% !TeX program = xelatex
\documentclass[sigconf, screen, review, anonymous]{acmart}

%% Chinese body text (XeLaTeX). Fall back to fontset=fandol if Windows fonts are unavailable.
\usepackage[UTF8,scheme=plain,fontset=windows]{ctex}

\usepackage{amsmath,amssymb}
\usepackage{booktabs}
\usepackage{multirow}

%% Rights / conference metadata (placeholders; internal draft).
\setcopyright{none}
\copyrightyear{2026}
\acmYear{2026}
\acmDOI{10.1145/XXXXXXX.XXXXXXX}
\acmConference[MM '26]{ACM Multimedia 2026}{TBD}{TBD}
\acmISBN{978-1-4503-XXXX-X/2026/XX}

\AtBeginDocument{%
  \providecommand\BibTeX{{Bib\TeX}}}

\begin{document}

\title[CHORD]{CHORD：面向多模态大语言模型的物体共振解码幻觉校准}
\subtitle{\textit{Calibrating Hallucinations via Object-Resonant Decoding in Multimodal Large Language Models}}

\author{占位作者}
\affiliation{%
  \institution{Anonymized for review}
  \country{USA}}

\begin{abstract}
多模态大语言模型（MLLM）在跨模态理解与生成上已取得显著进展，但仍持续受幻觉问题困扰，即生成与视觉输入不一致的文本。我们指出，其根本原因之一在于自回归解码阶段的\textbf{跨模态粒度失配}：图像被编码为连续的整体特征，而文本必须以离散、且常呈语义碎片化的子词（sub-word）token 序列输出；这使得难以将无独立语义的单 token 可靠地锚定到具体视觉实体，模型易进入未经验证的解码状态，语言先验的统计惯性进而放大幻觉。现有缓解策略往往依赖昂贵重训或高代价的多轮验证。为此，我们提出 \textbf{CHORD}（Calibrating Hallucinations via Object-Resonant Decoding，通过物体共振解码校准幻觉）——一种\textbf{免训练}、可即插即用的解码框架，借助 \textbf{Object-Resonant Decoding（物体共振解码）} 将生成轨迹重新锚定到视觉事实。CHORD 在活跃解码前沿上工作，动态评估\textbf{候选文本片段}与\textbf{外部视觉证据}的对齐程度：仅当片段与具体视觉物体在语义上``共振''时增强其概率，而对仅由语言惯性驱动的无根据内容予以抑制。为在控制开销的同时捕获\textbf{过度自信型}幻觉，我们采用\textbf{连续 Top-$1$ 监控}并在必要时\textbf{回退到全 $K$ 候选}的共振计算，并利用外部文本编码器的 KV 缓存复用，使完整前沿评分仅在 Top-$1$ 假设未通过共振检验时触发。在多个面向幻觉的基准上，CHORD 显著降低物体幻觉；作为与宿主架构解耦的框架，它无需额外训练，为可信 MLLM 部署提供可扩展路径。

\end{abstract}

\begin{CCSXML}
<ccs2012>
 <concept>
  <concept_id>10010147.10010178.10010224</concept_id>
  <concept_desc>Computing methodologies~Artificial intelligence</concept_desc>
  <concept_significance>500</concept_significance>
 </concept>
 <concept>
  <concept_id>10010147.10010178</concept_id>
  <concept_desc>Computing methodologies~Machine learning</concept_desc>
  <concept_significance>300</concept_significance>
 </concept>
</ccs2012>
\end{CCSXML}

\ccsdesc[500]{Computing methodologies~Artificial intelligence}
\ccsdesc[300]{Computing methodologies~Machine learning}

\keywords{多模态大语言模型, 幻觉缓解, 解码时干预, 物体共振解码, CHORD, 接地}

\maketitle

\section{引言}
\label{sec:intro}

多模态大语言模型（MLLM）已推动视觉问答、图像描述与开放式多模态推理快速发展，但即使是强模型仍易出现幻觉：生成流畅却接地不足、虚构图像中不存在的物体，或在强语言先验下对视觉不支持的实体过度自信。当正确证据稀疏、局部化或被无关视觉上下文淹没时，这类失败尤为有害。

近期工作表明，\textbf{推理期干预}是替代重训宿主模型的有效途径。对比解码（如 VCD）、基于注意力的控制（如 OPERA）以及层间对比解码（如 DoLa）均将自回归解码视为可操作的干预面。然而，仍存在更严格的问题：

\begin{quote}
\textbf{解码期缓解能否在异构 MLLM 家族间保持可迁移性——无需按模型单独训练，且不依赖宿主内部的隐状态、logit 或注意力？}
\end{quote}

该目标不同于单点提升骨干性能。任何消费模型特异多模态激活、并在宿主原生几何中写入校正的方法，即便视觉—语言权重冻结，仍属\textbf{宿主耦合}。\textbf{隐状态校准}类路线说明了这一现象，但难以满足严格的即插即用可迁移要求。

因此我们提出\textbf{可迁移性边界}：机制若完全运行于模型原生的隐空间、logit 空间或注意力空间内，则无法实现严格意义上的热插拔。可检验的替代方案是将所有评分限制在由\textbf{冻结的公开视觉特征}、\textbf{冻结文本编码器}、\textbf{解码得到的候选片段}以及\textbf{确定性}分词器桥接所构成的\textbf{共享外部空间}中，且\textbf{不}引入绑定于特定宿主的可学习适配器。

\textbf{CHORD}（Calibrating Hallucinations via Object-Resonant Decoding）体现该原则。它执行\textbf{候选驱动的视觉共振}：关键信号并非未完成前缀与图像的全局相似度，而是\textbf{当前候选片段}是否激活一致的局部视觉证据，以及该证据是否支持干预。缓解被刻画为对与候选绑定的外部、物体级证据的选择性调用，而非对无关区域或纯语言延续偏置的整体重分配。

由此有两点推论：其一，除原始幻觉降幅外，还应在\textbf{外部编码器与区域库固定、仅分词器桥接随宿主确定性变化}的前提下，考察 CHORD 在\textbf{不同宿主家族}上是否\textbf{稳定}；其二，基于稠密区域库的片段级共振应明显优于朴素的全局图像先验，否则额外机制难以自证必要。

本文贡献如下：
\begin{enumerate}
  \item 形式化限制宿主内部解码器严格可移植控制的\textbf{可迁移性边界}。
  \item 提出\textbf{CHORD}：\textbf{免训练}解码框架，在外部嵌入空间中利用物体共振证据与闭式\textbf{共振增量 $\Delta S$}校准活跃 logit 前沿。
  \item 引入\textbf{无参数分词器桥接}及\textbf{连续 Top-$1$ 监控}（配合全 $K$ 回退），使共振检验轻量，并在不依赖朴素熵门控的前提下抑制过度自信幻觉。
  \item 在 POPE、CHAIR 上给出含解码基线、消融与效率分析的\textbf{评估协议}。
\end{enumerate}

全文结构：第~\ref{sec:related} 节为相关工作；第~\ref{sec:method} 节给出边界与 CHORD 方法；第~\ref{sec:exp} 节为实验；第~\ref{sec:conclusion} 节为结论。

\section{相关工作}
\label{sec:related}

\subsection{多模态幻觉评估与基准}

视觉—语言模型中的幻觉评估已形成若干互补协议。\textbf{CHAIR}~\cite{rohrbach2018chair} 衡量生成描述中提及的物体是否被图像所支持，表明流畅文本仍可能视觉不成立。\textbf{POPE}~\cite{li2023pope} 将评估转化为平衡的物体存在性是非题，在受控提示下度量幻觉并揭示对语言先验的敏感。\textbf{AMBER}~\cite{amber2023} 将视角扩展到属性、关系等多类错误，避免将``仅物体级''分数等同于完整多模态忠实度。三者回答的问题不同：CHAIR 面向开放式描述；POPE 面向判别式探测；AMBER 强调多面失效模式，不宜互相替代。

更大规模的判别套件在任务与上下文结构上进一步扩展。\textbf{PhD}~\cite{liu2025phd} 提供大规模 VQA 式基准，含多种识别任务与对抗式提问模式，用于检验 MLLM 在视觉误导或先验触发场景下的鲁棒性。本文以 CHAIR 与 POPE 为主要汇报轴（第~\ref{sec:exp} 节），并将此类更广评估脉络作为未来压力测试的背景。

\subsection{解码期缓解与控制}

一类重要方法在\textbf{不更新权重}的前提下于自回归解码过程中干预。\textbf{视觉对比解码（VCD）}~\cite{leng2024vcd} 对比完整视觉与扰动视觉条件下的 logit，以抑制语言先验驱动的物体。\textbf{OPERA}~\cite{huang2024opera} 利用注意力与回溯信号约束生成轨迹。\textbf{DoLa}~\cite{chuang2024dola} 源于纯文本 LM，通过对比浅层与深层 logit 提升事实性；迁移到 MLLM 后仍是强的\textbf{内部}层间基线，但本身并不在外部视觉空间中强制物体级接地。

近期面向 MLLM 的解码改进仍主要\textbf{读取或重塑内部表示}。例如 \textbf{DeCo}~\cite{wang2024deco} 在浅层似乎仍保留视觉证据而深层被语言先验压制时，动态融合更早层的 logit。这类设计说明栈内仍存在有用视觉信号，但也继承了需访问宿主层输出与几何的可迁移性限制。

\textbf{CHORD} 与第~\ref{sec:method} 节强调的两条轴线一致：（i）\textbf{不}消费宿主隐状态或注意力张量——仅使用短候选前沿上的 logit 与确定性分词器前瞻；（ii）在\textbf{冻结}的区域级视觉库与\textbf{冻结}的公开文本编码器上对\textbf{候选片段}打分，使干预信号定义于共享外部空间。其针对子词提议与物体级证据之间的\textbf{粒度}鸿沟，并以连续 Top-$1$ 监控控制片段级评分开销。

\subsection{后处理、偏好优化及与 CHORD 的关系}

互补路线在生成\textbf{之后}修改输出，或直接更新模型。\textbf{LURE}~\cite{zhou2024lure} 基于与物体幻觉相关的统计因素做轻量后修。\textbf{Woodpecker}~\cite{yin2024woodpecker} 通过多阶段工具流水线抽取主张、查询外部视觉模块并重写不一致文本。\textbf{Volcano}~\cite{lee2024volcano} 迭代自反馈与修订。\textbf{HA-DPO}~\cite{zhao2023hadpo} 通过偏好优化持久改变模型行为。正因不受逐步解码预算约束，它们才能使用更丰富验证或更强监督。

\textbf{CHORD} 针对更窄设定：在 top-$K$ 前沿上\textbf{免训练}、\textbf{解码内}校准，依赖外部物体级证据与提升后片段上的闭式共振增量 $\Delta S$。时间轴上与解码基线~\cite{leng2024vcd,huang2024opera,chuang2024dola,wang2024deco} 最接近，证据形态上接近工具增强路线~\cite{yin2024woodpecker}，但避免完整后处理流水线且不做偏好更新~\cite{zhao2023hadpo}。第~\ref{sec:method} 节进一步说明该设计与可迁移性边界及抑制过度自信幻觉（不唯熵门控）的关系。

\section{方法}
\label{sec:method}

\subsection{可迁移性边界与粒度失配}

若解码期干预模块的输入或输出具有宿主特异性，则难以称为严格通用。早期校准方法读取模型原生隐状态并投影到宿主词表几何中，此类可学习映射按构造即与宿主耦合，无法在无重训情况下迁移到异构 MLLM。

要实现真正的热插拔，干预须运行于标准化的外部语义空间。设计外部控制器又必须直面幻觉的根本动因之一：\textbf{跨模态粒度失配}。MLLM 将图像编码为连续、整体化的特征图，而自回归生成却输出离散、常呈碎片化的子词 token（例如 \texttt{["gi", "\#\#raf", "\#\#fe"]}）。在连续视觉实体与无语义碎片 token 之间建立可靠物理接地极为困难。粒度崩塌使模型陷入\emph{开环解码}：视觉校验失效，语言先验惯性压倒视觉约束。CHORD 旨在弥合该失配并闭合解码回路。

\subsection{CHORD 框架概述}

CHORD（Calibrating Hallucinations via Object-Resonant Decoding）为免训练、即插即用解码框架。它不覆盖全词表分布，仅作用于宿主在步 $t$ 的活跃解码前沿
\begin{equation}
  \mathcal{F}_t = \mathrm{Top}\text{-}K(L_t),
\end{equation}
其中 $L_t$ 为原始 logit。流程含四阶段：（1）分词器桥接；（2）连续共振监控；（3）假设驱动共振增量；（4）logit 校准。

生成前，CHORD 对输入图像 $I$ 做一次离线提取：冻结的外部视觉编码器构建稠密局部区域特征库 $\mathcal{Z}_R = \{z_{R_1}, z_{R_2}, \dots, z_{R_N}\}$，作为后续共振核验的静态``物理事实''缓存。

\subsection{无参数分词器桥接}

为克服粒度失配，CHORD 不直接对原始 token ID 评分。设 $B_m$ 为与宿主 $m$ 关联的\textbf{无参数}分词器桥接。对每个候选 $c_k \in \mathcal{F}_t$，桥接利用分词器前缀树做确定性前瞻，将 token 提升为语义完整的文本片段 $A_k$，得到候选片段集 $\mathcal{S}_t = \{A_1,\dots,A_K\}$。桥接不含可学习参数，是将模型特异的碎片化提议转为跨模型、可物理接地的语义命题的关键。

\subsection{假设驱动的共振增量}

CHORD 将解码建模为物理共振过程：将前缀 $p$ 与候选片段 $A_k$ 拼接视为指向区域库 $\mathcal{Z}_R$ 的``语义声纳''。为严格避免\emph{空间错配}（前缀与候选对齐到无关区域），共振增量锚定在\textbf{同一}物理区域上。首先为联合假设选取最相关区域：
\begin{equation}
  r^* = \arg\max_{r \in \mathcal{Z}_R} \cos\bigl(E_{\text{text}}(p \oplus A_k), z_r\bigr),
\end{equation}
其中 $E_{\text{text}}$ 为冻结公开文本编码器。再定义\textbf{共振增量 $\Delta S$}，刻画加入 $A_k$ 的纯物理效应：
\begin{equation}
  \Delta S(A_k) = \cos\bigl(E_{\text{text}}(p \oplus A_k), z_{r^*}\bigr) - \cos\bigl(E_{\text{text}}(p), z_{r^*}\bigr).
\end{equation}

该式起到\textbf{无启发式}判别作用：（1）$\Delta S>0$ 表示视觉支持或正确空间关系；（2）$\Delta S<0$ 表示与图像矛盾（如虚构物体）；（3）$\Delta S\approx 0$ 对应抽象词或语法词，CHORD 可安全忽略而无需脆弱 POS 标注。

\subsection{克服过度自信：连续监控与系统优化}

朴素熵门控难以捕获\textbf{过度自信幻觉}（低熵仍错）。CHORD 弃用纯熵门控，采用\textbf{连续 Top-$1$ 监控与全 $K$ 回退}，并结合 $E_{\text{text}}$ 的 KV 缓存。前缀 $p$ 持续缓存；$p\oplus A_k$ 往往只需再前向 1--2 个 token。每步对宿主最自信 token 计算 $\Delta S(A_{\text{top1}})$：若 $\Delta S(A_{\text{top1}})\ge 0$ 则直接接受；若 $<0$ 则触发对全部 $K$ 个候选的完整共振评估以寻找接地替代。该策略在抑制过度自信的同时控制计算开销。

\subsection{Logit 校准}

触发完整干预时，CHORD 将共振证据注入宿主 logit。与片段 $A_k$ 对应的候选 token $c_k$ 的调整 logit 为
\begin{equation}
  \tilde{L}_t(c_k) = L_t(c_k) + \alpha \cdot \Delta S(A_k),
\end{equation}
$\alpha$ 为强度超参；再经 softmax 归一化。有界重分配使轨迹向视觉事实偏转，同时尽量不破坏宿主语言结构且无需更新参数。

\section{实验}
\label{sec:exp}

为验证 CHORD，实验围绕四个研究问题展开：

\textbf{RQ1（有效性与可迁移性）：} CHORD 能否在\textbf{无模型特异重训}的前提下，在异构 MLLM 家族上一致缓解物体幻觉？

\textbf{RQ2（泛化性）：} CHORD 能否推广到开放式生成，而不损害宿主的结构与语言能力？

\textbf{RQ3（机制消融）：} 分词器桥接、共振增量 $\Delta S$ 等各组件如何贡献最终性能？

\textbf{RQ4（系统效率）：} 连续 Top-$1$ 监控能否将额外计算控制在可忽略水平？

\subsection{实验设置}

\textbf{模型与基线。} 我们在 Qwen-VL（如 Qwen-VL-Chat）、LLaVA-1.5、InstructBLIP 三类家族上评估。对比基线包括标准自回归解码（Base）及三种免训练解码基线：\textbf{DoLa}、\textbf{OPERA}、\textbf{VCD}。

\textbf{基准。} 使用 \textbf{POPE}（random / popular / adversarial 划分）做系统物体幻觉评估，\textbf{CHAIR} 评估开放式图像描述中的幻觉率。

\textbf{实现细节。} CHORD 无参数更新。区域缓存使用 Grounding DINO 与冻结 CLIP ViT-L/14；共振探针使用 CLIP 文本编码器。

\subsection{主结果：幻觉缓解（RQ1）}

表~\ref{tab:pope} 汇总 POPE（random 划分）主结果。CHORD 相对 Base 一致提升；在外部、与宿主内部几何解耦的语义空间中计算 $\Delta S$，可抑制语言先验驱动错误。POPE 的 popular 与 adversarial 划分及扩展结果见补充材料，CHORD 保持稳定收益。

\subsection{泛化与结构保持（RQ2）}

仅在判别任务（如 POPE）上降幻觉不足；若破坏流畅描述能力则不可取。表~\ref{tab:chair} 报告 CHAIR-s（句级）与 CHAIR-i（实例级）。有界 logit 校准旨在仅惩罚与视觉矛盾实体，尽量保留宿主句法与推理习惯。

\subsection{消融实验（RQ3）}

表~\ref{tab:ablation} 在 Qwen-VL（POPE）上消融。\textbf{去掉分词器桥接、}直接对原始子词做共振会导致性能崩溃，说明弥合粒度失配是外部评估的前提。\textbf{用绝对相似度替代 $\Delta S$} 易产生空间错配（前缀与候选对齐到不同区域），F1 下降。\textbf{仅用全局图像特征、}不用稠密区域库则损失细粒度证据，缓解效果减弱。

\subsection{系统效率（RQ4）}

表~\ref{tab:efficiency} 在 Qwen-VL 上审计开销。每步全 $K$ 共振代价高；\textbf{连续 Top-$1$ 监控}在 Top-$1$ 假设通过检验时以 $O(1)$ 过滤安全或非视觉 token，仅在未通过时回退全 $K$。结合文本编码器 KV 缓存，每步额外延迟相对宿主 MLLM 总体推理保持较小。

\begin{table*}[t]
  \centering
  \caption{POPE 基准（random 划分）主结果。TBD 表示定稿待填数值。}
  \label{tab:pope}
  \small
  \begin{tabular}{llcccc}
    \toprule
    宿主模型 & 方法 & Acc & Prec & Recall & F1 \\
    \midrule
    \multirow{5}{*}{\textbf{Qwen-VL}}
      & Base & TBD & TBD & TBD & TBD \\
      & DoLa & TBD & TBD & TBD & TBD \\
      & OPERA & TBD & TBD & TBD & TBD \\
      & VCD & TBD & TBD & TBD & TBD \\
      & \textbf{CHORD（本文）} & \textbf{0.9153} & \textbf{0.9822} & \textbf{0.8460} & \textbf{0.9090} \\
    \midrule
    \multirow{5}{*}{\textbf{LLaVA-1.5}}
      & Base & TBD & TBD & TBD & TBD \\
      & DoLa & TBD & TBD & TBD & TBD \\
      & OPERA & TBD & TBD & TBD & TBD \\
      & VCD & TBD & TBD & TBD & TBD \\
      & \textbf{CHORD（本文）} & \textbf{TBD} & \textbf{TBD} & \textbf{TBD} & \textbf{TBD} \\
    \midrule
    \multirow{2}{*}{\textbf{InstructBLIP}}
      & Base & TBD & TBD & TBD & TBD \\
      & \textbf{CHORD（本文）} & \textbf{TBD} & \textbf{TBD} & \textbf{TBD} & \textbf{TBD} \\
    \bottomrule
  \end{tabular}
\end{table*}

\begin{table}[t]
  \centering
  \caption{CHAIR 开放式生成评估（越低越好）。}
  \label{tab:chair}
  \small
  \begin{tabular}{llcc}
    \toprule
    宿主模型 & 方法 & CHAIR-s $\downarrow$ & CHAIR-i $\downarrow$ \\
    \midrule
    \multirow{2}{*}{Qwen-VL} & Base & TBD & TBD \\
      & \textbf{CHORD} & \textbf{TBD} & \textbf{TBD} \\
    \multirow{2}{*}{LLaVA-1.5} & Base & TBD & TBD \\
      & \textbf{CHORD} & \textbf{TBD} & \textbf{TBD} \\
    \bottomrule
  \end{tabular}
\end{table}

\begin{table}[t]
  \centering
  \caption{Qwen-VL 上 CHORD 组件消融（POPE）。}
  \label{tab:ablation}
  \small
  \begin{tabular}{lcc}
    \toprule
    变体 & F1 & $\Delta$ F1 \\
    \midrule
    \textbf{完整 CHORD} & \textbf{TBD} & -- \\
    无分词器桥接（原始 token） & TBD & TBD \\
    绝对 $S_{\text{new}}$（无 $\Delta S$） & TBD & TBD \\
    仅全局特征（无区域库） & TBD & TBD \\
    \bottomrule
  \end{tabular}
\end{table}

\begin{table}[t]
  \centering
  \caption{Qwen-VL 推理效率审计。}
  \label{tab:efficiency}
  \small
  \begin{tabular}{lccc}
    \toprule
    方法 & 峰值显存 (GB) & 每步额外延迟 (ms) & F1 \\
    \midrule
    Base（束搜索） & TBD & 0 & TBD \\
    DoLa & TBD & TBD & TBD \\
    CHORD（每步全 $K$） & TBD & TBD（高） & TBD \\
    \textbf{CHORD（Top-$1$ 监控）} & \textbf{TBD} & \textbf{TBD} & \textbf{TBD} \\
    \bottomrule
  \end{tabular}
\end{table}

\begin{figure}[t]
  \centering
  \fbox{%
    \begin{minipage}{0.92\linewidth}
      \centering
      \vspace{0.8em}
      \textbf{（插图占位）}CHORD 流水线：宿主 Top-$K$ 前沿 $\rightarrow$ 无参数分词器桥接 $\rightarrow$ 冻结文本/视觉编码器与区域库 $\rightarrow$ 共振增量 $\Delta S$ $\rightarrow$ 对 $\mathcal{F}_t$ 的有界 logit 校准。
      \vspace{0.8em}
    \end{minipage}}
  \caption{解码期 CHORD 接入示意图（占位）。}
  \label{fig:overview}
\end{figure}

\section{结论}
\label{sec:conclusion}

本文指出现代 MLLM 的一个结构性问题：\textbf{跨模态粒度失配}。连续视觉特征与离散、碎片化的子词 token 之间的不一致，使自回归生成易处于未经验证的\emph{开环解码}状态，语言先验惯性常压倒视觉约束，从而催化物体与关系幻觉。

针对该根因，我们提出 \textbf{CHORD}（Calibrating Hallucinations via Object-Resonant Decoding）：通过无参数分词器桥接将碎片 token 提升为语义完整片段，在外部共享语义空间中计算\textbf{共振增量 $\Delta S$}，在避免脆弱启发式与沉重黑盒评估器的同时，强化视觉一致内容并抑制虚构描述。

我们进一步用配备 KV 缓存加速的\textbf{连续 Top-$1$ 监控}替代朴素熵门控，在控制开销的同时仍能有效应对过度自信型幻觉。

多宿主家族上的实验表明，CHORD 在显著抑制物体幻觉的同时，尽量保持语言流畅性。结果表明，通过外部物体共振解码可实现强且可跨模型的接地控制，为更可信的 MLLM 部署提供可扩展范式。


\appendix
\section{补充材料}

POPE 在 popular 与 adversarial 划分上的补充结果、实现细节与超参设置见补充材料。正文表~\ref{tab:pope} 报告 random 划分；CHORD 在其他划分上保持一致的收益趋势。


\bibliographystyle{ACM-Reference-Format}
\bibliography{references}

\end{document}
